\section*{Limitations}

% Style note (v2, 2026-05-07): rewritten in the EMNLP 2025 Outstanding-paper
% rhetorical mode --- limitations as competing interpretations and scope
% caps, not failures. Each paragraph either preempts a known reviewer
% concern (efficiency-axis choice, BM25 first-stage, English-only data) or
% surfaces a deliberate scope decision (MaxContext-only Phase F, pool
% ceiling, generation-only scoring). Word target ~250 to fit the
% Limitations budget per plan v4.

Our central efficiency claim is restricted to serial LLM calls, and our central effectiveness claim to per-query nDCG@10 under a specific non-inferiority margin. This is a deliberate narrowing. One reading of a positive result would be that whole-pool DualEnd is uniformly more efficient than windowed setwise reranking; another, equally reasonable, is that the call-count gap matters most when serial latency dominates and is largely absorbed by batching, prompt re-encoding, or token throughput. We report wall-clock and prompt/completion/total tokens alongside calls so that either reading is auditable, but our findings alone cannot adjudicate which efficiency axis matters for a given deployment.

We rerank only the BM25 top-$N$ pool. This isolates the reranker's contribution given a strong but fixed first stage, but it also means that any per-query effect we observe is a property of \emph{both} the reranker and the BM25 pool's composition. Whether the same pattern holds for dense first-stage retrievers, hybrid sparse--dense pools, or weaker first stages is not a question this paper can answer.

We restrict MaxContext to generation-based scoring on causal LLMs from three open-weight families, on TREC DL19/20 and six BEIR domains, all in English. The joint best--worst objective in DualEnd is not equivalent to an independent likelihood score, so we do not include T5-style likelihood scorers in the matrix. Scaling MaxContext to non-English benchmarks, mixed-multilingual pools, and likelihood- or embedding-based scorers is a question we leave open.

Pool sizes are capped at $N{=}100$ and gated on a per-family context-fit preflight; the cap is context-window- and passage-length-dependent rather than fundamental. Methodologically, the position-bias controls (\S\ref{sec:results-position-bias}) are applied only to MaxContext, not to the windowed setwise baselines, because their per-comparison semantics differ. We report the deterministic 10$\times$ stability check and the Phase F forward/reverse/shuffle deltas as separate views of input-order behavior, and frame conclusions about position-bias under whole-pool prompting accordingly.
